
\def\PUDcodigo{ECA202}

\def\PUDcodacad{04507.40}

\def\PUDdisciplina{Controle II}

\def\PUDsemestre{S8}

\def\PUDhoras{80}

%NOVOS ---------------------------------------------------
\def\PUDhorasTeoria{70}
\def\PUDhorasPratica{10}

\def\PUDinterECA{ 
\hyperref[PUD:ECA213]{Microcontroladores (04507.28)}

\hyperref[PUD:ECA233]{Controle I (04507.31)}

\hyperref[PUD:ECA205]{PDS (04507.34)}

\hyperref[PUD:ECA235]{Controle III (04507.58)} 
}

\def\PUDatuaisECA{---}

\def\PUDrecursos{Projetor multimídia; Quadro branco e pincel; Aulas em laboratório.}

%----------------------------------------------------------

\def\PUDcreditos{4}

\def\PUDprerequisito{\hyperref[PUD:ECA233]{Controle I (04507.31)}}

\def\PUDpreacad{\hyperref[PUD:ECA233]{Controle I (04507.31)}}

\def\PUDobjetivos{Compreender as ferramentas básicas de análise e projeto de sistemas de controle digital.  Aplicar as ferramentas na resolução de problemas afins.}

\def\PUDementa{Breve revisão do domínio da frequência. Introdução ao controle digital.  Breve revisão de princípios de controle e de análise de sinais e de sistemas discretos.  Sistemas amostrados.  Equivalentes discretos.  Sistemas de tempo discreto.  Transformada Z modificada.  Resposta temporal e sistemas discretos.  Estabilidade.  Projeto de controladores digitais.  Controle ótimo linear-quadrático.  Efeitos de quantização.  Hierarquia de sistemas de controle.  Estratégias de controle.  Implantação de sistemas de controle e automação industrial.  Critérios de desempenho, caracterização e sintonia de controladores industriais.  Simulações e realização de práticas envolvendo metodologias de projetos de controladores discretos.}

\def\PUDconteudo{ & Análise no domínio da frequência. Diagramas de Bode; gráficos polares; Critérios de estabilidade de Nyquist. Resposta em frequencia em malha aberta e em malha fechada.
\\
& Controle Analógico Versus Controle Digital;  Sistemas Típicos de Controle Digital;  Definições;  Quantização: Aquisição e Conversão de Sinal Digital para Analógico;  Exemplos de Sistemas Controlados: Sistemas de Controle Monovariáveis e Sistemas de Controle Multivariáveis. 
\\ 
 & Análise de Sistemas de Controle Discreto: Funções de Transferência: Função de Transferência do Hold, Função Simples, Elementos em Cascata, Malha Fechada e Controlador Digital;  Resposta Transitória e de Estado Permanente: Especificações de Resposta Transitória ao Degrau, Mapeamento entre Planos s e Plano z, Análise de Erro em Estado Permanente, Efeito de Perturbação na Planta;  Realização de Controladores Digitais: Programação Direta, Programação Padrão. 
\\ 
 & Controladores PID discretos, PID no formato RST, versões aprimoradas do controlador, exemplos de projeto e sintonia, aplicação em modelos de plantas Industriais.
 %Projeto de Controladores Digitais por Métodos Convencionais:  Efeito das Ações de Controle;  Digitalização de Controladores Analógicos: Aproximação Numérica da Integração e Aproximação Numérica da Diferenciação;  Filtragem da Entrada Analógica da Planta;  Estabilidade de Sistemas Controlados: Localização de Pólos e Estabilidade, Teste de Estabilidade de Jury, Critério de Estabilidade de Routh;  Lugar das Raízes. 
\\ 
 & Análise de Sistemas de Controle no Espaço de Estados: Controlabilidade e Observabilidade. 
\\ 
 & Projeto de Sistemas de Controle no Espaço de Estados: Alocação de pólos;  Observadores de estado;  Projeto de sistemas reguladores com observadores;  Projeto de sistemas de controle com observadores. }

\def\PUDmetodo{Aulas expositórias que podem ser teóricas e/ou práticas, onde as práticas no laboratório serão marcadas no decorrer da disciplina.}

\def\PUDavalia{A avaliação será feita com aplicação de provas teóricas e/ou práticas, além da possibilidade de inclusão de trabalhos, elaboração de artigos técnicos e seminários no decorrer da disciplina.
%A avaliação será feita com aplicação de provas teóricas e/ou práticas, além da possibilidade de inclusão de trabalhos e seminários no decorrer da disciplina.
}

\def\PUDbibbasica{ & \href{www.biblioteca.ifce.edu.br/index.asp?codigo_sophia=63123}{Nise}, Norman S.  Engenharia de sistemas de controle.  6º ed.  Rio de Janeiro, RJ: LTC, 2012.  745p.  ISBN: 9788521621355.
\\
& \href{www.biblioteca.ifce.edu.br/index.asp?codigo_sophia=62612}{Franklin}, Gene F. Sistemas de controle para engenharia.  6º ed.  Porto Alegre, RS: Bookman, 2013.  702p.  ISBN: 9788582600672.
\\
& \href{www.biblioteca.ifce.edu.br/index.asp?codigo_sophia=63228}{Ogata}, Katsuhiko.  Engenharia de controle moderno.  5º ed.  São Paulo, SP: Pearson Prentice Hall, 2010.  809p.  ISBN: 9788576058106. 
 }

\def\PUDbibcomplementar{ & \href{www.biblioteca.ifce.edu.br/index.asp?codigo_sophia=40232}{Oppenheim}, Alan V.; Schafer, Ronald W.  Processamento em tempo discreto de sinais.  3ª ed.  São Paulo, SP: Pearson Education do Brasil, 2012.  665p.  ISBN: 9788581431024.  
%&  DORF, Richard C.;  BISHOP, Robert H.  Sistemas de controle modernos.  12. ed.  Rio de Janeiro (RJ): LTC, 2013.  659 p.  ISBN 9788521619956
\\
 &  \href{www.biblioteca.ifce.edu.br/index.asp?codigo_sophia=62576}{Leonardi}, F.; Maya, Paulo A.  Controle essencial.  2º ed.  São Paulo, SP: Pearson Education do Brasil, 2014.  347p.  ISBN: 9788543002415.
\\
 & \href{www.biblioteca.ifce.edu.br/index.asp?codigo_sophia=39837}{Lathi}, B. P.  Sinais e Sistemas Lineares.  2º ed.  Porto Alegre, RS: Bookman, 2007.  856p.  ISBN: 9788560031139.
%\\ 
 %&   P.  E.  Wellstead, M.  B.  Zarrop;  "Self-tuning Systems: Control and Processign Signal".  Wiley, 1991.  ISBN 9780471928836. }
\\
 & \href{www.biblioteca.ifce.edu.br/index.asp?codigo_sophia=56037}{Shumway-Cook}, Anne; Woollacott, Marjorie H.  Controle Motor: teoria e aplicações práticas.  E-book.  Manole.  636p.  ISBN: 9788520427477.
\\
 & \href{www.biblioteca.ifce.edu.br/index.asp?codigo_sophia=55618}{Gimenez}, Salvador P.  Microcontroladores 8051: teoria do Hardware e do Software: aplicações em controle digital: laboratório e simulação.  E-book.  Pearson.  272p.  ISBN: 9788587918284. }

\def\PUDautor{Geraldo Luis Bezerra Ramalho}

\def\PUDdata{2013-05-21}

\def\PUDrevisao{3}

\def\PUDrevisor{Luiz Daniel Santos Bezerra}

\def\PUDdatarevisao{2019-05-15}

\PUDcorpo
